% SEGMENT OLP-0374-S01
% Part: lambda-calculus
% Chapter: lambda-definability
% Section: introduction

\documentclass[../../../include/open-logic-section]{subfiles}

\begin{document}

\olfileid{lam}{ldf}{int}
\olsection{அறிமுகம்}

முதல் பார்வையில், லாம்டா கலனம் என்பது சார்புகளைக் குறிக்கும் கோவைகளையும்,
அவற்றை மற்றவற்றின்மீது செயல்படுத்துவதையும் பற்றிய மிக அருவமான ஒரு கலனம்
மட்டுமே. இவை குறிப்பிட்ட வகைப் பொருள்களின் சார்புகள் என்று லாம்டா கலனத்தின்
தொடரமைப்பில் எதுவும் உணர்த்துவதில்லை; குறிப்பாக, தொடரமைப்பில் இயல் எண்கள்
பற்றிய குறிப்பே இல்லை. அதன் அடிப்படைச் செயல்களான செயல்படுத்தலும் லாம்டா
அருவமாக்கமும் இயல் எண்கள்மீதான சார்புகளுக்கு மட்டுமல்லாமல் எந்தச்
சார்புக்கும் பொருந்தும் செயல்களாகும்.

எனினும், சிறிது புத்திக்கூர்மையுடன் லாம்டா கலனத்தில் எண்கணிதச் சார்புகளை,
அதாவது இயல் எண்கள்மீதான சார்புகளை, வரையறுக்க முடியும். இதற்காக, ஒவ்வோர் இயல்
எண்~$n \in \Nat$ க்கும் $n$ இன் \emph{சர்ச் எண்ணுரு} எனப்படும் ஒரு சிறப்பு
$\lambd$-சொல்~$\num n$ ஐ வரையறுக்கிறோம். (சர்ச் எண்ணுருக்கள் அலான்சோ சர்ச்சின்
பெயரால் அழைக்கப்படுகின்றன.)

\begin{defn}
  $n \in \Nat$ எனில், அதற்குரிய \emph{சர்ச் எண்ணுரு} $\num{n}$ ஆனது~$n$ ஐக்
  குறிக்கிறது:
  \[
    \num{n} \ident \lambd[fx][f^n(x)]
  \]
  இங்கு $f^n(x)$ என்பது $f$ ஐ~$x$ மீது $n$ முறை செயல்படுத்துவதன் விளைவைக்
  குறிக்கிறது. எடுத்துக்காட்டாக, $\num{0}$ என்பது $\lambd[fx][x]$; மேலும்
  $\num{3}$ என்பது $\lambd[fx][f(f(f\,x))]$.
\end{defn}

சர்ச் எண்ணுரு $\num n$, இரண்டு உள்ளீடுகள் $f$, $x$ ஐ ஏற்று $f^n(x)$ ஐத்
திருப்பித்தரும் சார்பைக் குறிக்கும் ஒரு லாம்டா சொல்லாகக் குறியாக்கப்பட்டுள்ளது.
சர்ச் எண்ணுருக்கள் தெளிவாக இயல்புவடிவத்தில் உள்ளன.

லாம்டா கலனத்தில் இயல் எண்களுக்கான பிரதிநிதித்துவம், அவற்றைக்கொண்டு கணக்கிட
முடிந்தால் மட்டுமே பயனுள்ளது. லாம்டா கலனத்தில் சர்ச் எண்ணுருக்களைக்கொண்டு
கணக்கிடுவது என்பது, ஒரு $\lambd$-சொல்~$F$ ஐ அத்தகைய சர்ச் எண்ணுருவின்மீது
செயல்படுத்தி, ஒருங்கிணைந்த சொல்~$F\, \num n$ ஐ ஓர் இயல்புவடிவத்திற்குக்
குறைப்பதாகும். அது எப்போதும் ஓர் இயல்புவடிவத்திற்குக் குறைந்து, அந்த
இயல்புவடிவம் எப்போதும் ஒரு சர்ச் எண்ணுரு~$\num m$ ஆக இருந்தால், கணக்கீட்டின்
வெளியீட்டை எண்~$m$ என்று கருதலாம். பின்னர்~$F$ ஆனது ஒரு சார்பு
$f\colon \Nat \to \Nat$ ஐ வரையறுக்கிறது என்று கருதலாம்; துல்லியமாக,
$F\, \num n \red \num m$ எனில் மட்டுமே $f(n) = m$ ஆக அமையும் சார்பு.
சர்ச்--ரோசர் பண்பினால், இயல்புவடிவங்கள் இருந்தால் அவை தனித்தவை. ஆகவே
$F\, \num n \red \num m$ எனில், $F \, \num n$ வேறு எந்த
இயல்புவடிவத்திற்கும், குறிப்பாக வேறு எந்தச் சர்ச் எண்ணுருவிற்கும், குறைய
முடியாது.

மறுதிசையில், ஒரு சார்பு $f\colon \Nat \to \Nat$ தரப்பட்டால், இவ்விதத்தில்~$f$
ஐ வரையறுக்கும் சொல் $F$ ஒன்று உள்ளதா என்று கேட்கலாம். அந்நிலையில் $F$ ஆனது~$f$
ஐ \emph{!!{lambda define}s} என்றும், $f$ ஆனது !!{lambda definable} என்றும்
கூறுகிறோம். இதைப் பல-இடச் சார்புகளுக்கும் பகுதிச் சார்புகளுக்கும்
பொதுமைப்படுத்தலாம்.

\begin{defn}
% SOURCE CORRECTION TA-LAM-033: the undefined branch makes f partial, so use the edition's partial-function arrow rather than a total-function arrow.
$f\colon \Nat^k \pto \Nat$ என்று கொள்க. எல்லா
$n_0$, \dots,~$n_{k-1}$ க்கும் பின்வருவன அமையுமானால், ஒரு லாம்டா சொல்~$F$
ஆனது $f$ ஐ \emph{!!{lambda define}s} என்கிறோம்:
\[
F \, \num{n_0} \, \num{n_1} \dots \num{n_{k-1}} \red \num{f(n_0, n_1, \dots,
  n_{k-1})}
\]
$f(n_0, \dots, n_{k-1})$ வரையறுக்கப்பட்டால் மேலுள்ள குறைத்தல் அமையும்;
இல்லையெனில் $F \, \num{n_0} \,
\num{n_1} \dots \num{n_{k-1}}$ க்கு இயல்புவடிவம் இல்லை.
\end{defn}

மாறிலிச் சார்புகள் மிக எளிய ஓர் எடுத்துக்காட்டு. சொல் $C_k \ident
\lambd[x][\num{k}]$ ஆனது, $c_k(n) = k$ ஆக அமையும் சார்பு
$c_k\colon \Nat \to \Nat$ ஐ !!{lambda define}s.
% SOURCE CORRECTION TA-LAM-034: use c_k in the defining equation, matching the named constant function and its displayed type.
ஏனெனில் எந்த~$n$ க்கும் $C_k \, \num n \ident
(\lambd[x][\num{k}])\num n \redone \num{k}$. சமனிச் சார்பு
$\lambd[x][x]$ ஆல் !!{lambda defined} ஆக உள்ளது. மேலும் சிக்கலான சார்புகளை
வரையறுப்பது இயல்பாகவே கடினம்; பெரும்பாலும் மிகுந்த புத்திக்கூர்மை தேவைப்படும்.
எனவே ஒவ்வொரு கணிக்கத்தக்க சார்பும் !!{lambda definable} என்பது ஒருவேளை
வியப்பூட்டலாம். மறுதிசையும் மெய்: ஒரு சார்பு !!{lambda definable} எனில்,
அது கணிக்கத்தக்கது.

\end{document}
